\section{Sharp rank on the actual coexact domain}\label{sec:sharp-coexact}
The general-field obstruction in Section~\ref{sec:common-rank} counts full
row rank $hR$. On a coexact domain those rows have relations, so that count
cannot simply be reused. We calculate their actual annihilator instead.

\subsection{Closed fields and the annihilator formula}
For one internal channel, $\mathcal U_0^\perp=\ker d_1$ by finite
adjoint duality. We first prove the real decomposition
\begin{equation}\label{eq:closed-decomposition}
 \ker d_1=\{\text{constant vector fields}\}\oplus\operatorname{im}G,
 \qquad (G\phi)_j=D_j^+\phi.
\end{equation}
Complexify and use the finite characters
$\exp(2\pi\mathrm i k\cdot x/n)$. At frequency $k$ put
$q_j(k)=\exp(2\pi\mathrm i k_j/n)-1$. The closed equations are
$q_i\widehat u_j-q_j\widehat u_i=0$. If $k\ne0$, some $q_\ell\ne0$,
and the equations force $\widehat u_j=q_j\widehat u_\ell/q_\ell$ for
all $j$, including those with $q_j=0$. A coordinate-free choice of potential is
\[
 \widehat\phi(k)=
 \frac{\sum_j\overline{q_j(k)}\widehat u_j(k)}{\sum_j|q_j(k)|^2}
 \quad(k\ne0),\qquad \widehat\phi(0)=0.
\]
For real $u$ this formula has conjugate symmetry, hence gives a real
mean-zero $\phi$. Its gradient supplies every nonzero-frequency mode;
the zero mode is a constant vector. Conversely constants and gradients
are closed. A periodic gradient has zero component means, so its
intersection with constant vectors is zero. The kernel of $G$ consists
of constants because forward coordinate steps connect the torus.
This proves \eqref{eq:closed-decomposition}, including $n=1$.
It also yields the useful real dimension identity
\begin{equation}\label{eq:coexact-dimension}
 \dim\mathcal U_0=DV-[D+(V-1)]=(D-1)(V-1).
\end{equation}
Real and complex dimensions here agree by complexification of a real map;
there is no doubling from conjugate frequencies. Forward Nyquist symbols
are $-2$, rather than the spurious zero of a central difference.

For a partition into $R$ nonempty sets, let $S$ be the $R$-dimensional
scalar block-constant space and $W=S^D$. The adjoint range of its full
mean map is $W$: all row scaling factors in \eqref{eq:rows} are nonzero.
Define the real potential space
\[
 \mathcal F=\{\phi:\sum_x\phi(x)=0,\ G\phi\in W\}.
\]
Constants lie in $W$, and all gradient component means vanish. Therefore
\eqref{eq:closed-decomposition} gives the direct sum
\[
 W\cap\mathcal U_0^\perp
 =\{\text{constant vectors}\}\oplus G\mathcal F.
\]
Since $G$ is injective on mean-zero functions, this intersection has
dimension $D+\dim\mathcal F$. The adjoint of the restricted map is
$P_{\mathcal U_0}C^*$. Finite rank-nullity consequently gives the exact formula
\begin{equation}\label{eq:annihilator-rank}
 \rank(C|_{\mathcal U_0})=DR-D-\dim\mathcal F.
\end{equation}
This formula is for the actual restricted image, not its listed output size.

\subsection{Forward maximum principle and sharpness}
Let $\mathcal L=\sum_jD_j^+$. If a real $\phi$ satisfies
$\mathcal L\phi=0$, choose a site of global maximum. Every forward
increment there is nonpositive and their sum is zero, so each is zero.
Every forward neighbor is another maximum. Iterating these steps reaches
every torus site, including inverse movements by periodic wrap. Thus
$\phi$ is constant. At $n=1$ this is immediate. This is the classical
finite irreducible-chain maximum principle; see Levin--Peres,
Lemma~1.16~\cite{levinperes}. Here $\mathcal L=D(P_{\rm walk}-I)$ for
the uniform forward-coordinate walk. Neither reversibility nor
aperiodicity is needed. We claim no new maximum principle.

For $\phi\in\mathcal F$, every component of $G\phi$ lies in $S$.
Thus $\mathcal L\phi\in S$, with global sum zero by telescoping.
The maximum principle proves that
\[
 \mathcal L:\mathcal F\longrightarrow S\cap1^\perp
\]
is injective. The target has dimension $R-1$, since the global-sum
functional on $S$ is the nonzero weighted sum of block values with weights
$|B|$. Hence $\dim\mathcal F\le R-1$. Substitution into
\eqref{eq:annihilator-rank} proves $(D-1)(R-1)$ for one channel.
The input and output maps for $g$ independent internal channels are literal
direct sums, so ranks add and prove \eqref{eq:sharp-rank}.

For sharpness set $n=R$ and $B_a=\{x:x_1=a\}$.
Every $\phi(x)=f(x_1)$ with $\sum_af(a)=0$ belongs to $\mathcal F$:
its first forward derivative is constant on each fibre, and all others
vanish. These potentials form an $(R-1)$-dimensional space. The preceding
upper bound gives equality in \eqref{eq:annihilator-rank}, and then equality
in \eqref{eq:sharp-rank}. At $R=1$ the common global means vanish on the
coexact domain and both sides are zero. Sharpness concerns the allowed
family of tori with $n=R$, not every fixed pair $(n,R)$.

\subsection{A complete fixed-size coexact coefficient family}
For the coefficient consequence we retain $D\ge3$ and $m\ge4$.
Let $L=16$, $g=m$, $h=Dg$, $n=m(L+3)+3$ and $V=n^D$.
Use the compact product field in \eqref{eq:bumpfield}, with internal
unit vector $e_1\in\R^g$, translated by
\[
 q_t=(1+t(L+3),1,\ldots,1),\quad0\le t<m,
 \qquad B_t=q_t+\{0,\ldots,L\}^D.
\]
Lemma~\ref{lem:bump} and its preceding support proof apply to this exact
field: the norm is nonzero, the unit-gradient Rayleigh bound is $K_D/L^2$,
and every nonzero original forward start lies in $B_t$. The last
first-coordinate endpoint is $m(L+3)-2<n-1$, and the other endpoints are
$L+1<n-1$. Successive boxes have two unused sites between them.
Thus there are no wrap collisions and complete field and gradient-start
supports are disjoint. Their span $E_m\subset\mathcal U$ is truly
$m$-dimensional. Only spatial components $1,2$ of internal channel $1$
are active on $E_m$.

Put $r=Dh$, $M=rV$, $\sigma=\alpha_D^{1/P}c$ and define
\begin{equation}\label{eq:coexact-family}
 \begin{gathered}
 a_0=\min\{\sigma/2,cL^2/(4K_D)\}>0,\qquad H=\sigma M,\\
 A_{(i,j,\alpha),(i',j',\alpha')}(x)=
 \delta_{ii'}\delta_{jj'}\delta_{\alpha\alpha'}\lambda_{i,j,\alpha}(x),\\
 \lambda_{i,j,\alpha}(x)=
 \begin{cases}
 a_0,&x\in\bigcup_tB_t,\ j\in\{1,2\},\ \alpha=1,\\
 H,&\text{otherwise}.
 \end{cases}
 \end{gathered}
\end{equation}
The triple denotes gradient direction, spatial field component and
internal channel. This defines an actual positive diagonal matrix on all
$r=D^2g$ gradient coordinates at every start. The constant $K_D$ is the
bump Rayleigh constant of Lemma~\ref{lem:bump}, not the observation penalty.
Every nonzero derivative of each bump has coefficient $a_0$, so on $E_m$
\begin{equation}\label{eq:fixed-coexact-low}
 Q_A(U)\le a_0(K_D/L^2)\sum_x\|U(x)\|^2
                  \le(c/4)\sum_x\|U(x)\|^2.
\end{equation}
All norm and energy cross terms vanish by disjoint supports.

The number of weak local direction/component/site channels is exactly
$W_*=2Dm(L+1)^D=\kappa_Dm$, with $\kappa_D=2D17^D$.
This count includes zero-derivative entries assigned $a_0$; none are dropped.
The boxes are disjoint and $r\ge2D$, so $W_*\le M$.
Each weak cap cost is one; every strong cost is $M^{-P}$. Hence the full budget is
\begin{equation}\label{eq:fixed-coexact-budget}
 F_c(A)=W_*+(M-W_*)M^{-P}\in[\kappa_Dm,\kappa_Dm+1].
\end{equation}
The strong tail is bounded by $M^{1-P}\le1$, using $P>1$.
No arithmetic-mean or envelope condition is asserted for this matrix family.

If a common partition with $R$ members restores coercivity on $\mathcal U$
with $\Phi(0)=0$, its restriction to $E_m$ has rank at most $2R$.
Indeed every row has only the two active real component coordinates there.
If $2R<m$, a nonzero low-space kernel vector contradicts coercivity and
\eqref{eq:fixed-coexact-low}, after multiplying by $V$.
Thus $R\ge\lceil m/2\rceil$. The sharp rank bound independently yields
\[
 \rank(C_{\mathcal R}|_{\mathcal U})
 \ge g(D-1)(\lceil m/2\rceil-1)
 \ge g(D-1)m/4= h(D-1)m/(4D).
\]
Here $\lceil m/2\rceil-1\ge m/4$ for $m\ge4$.
Using $F_c(A)\le(\kappa_D+1)m$ proves \eqref{eq:coexact-product} and
completes Theorem~\ref{thm:sharp-coexact}. With $g=m$, necessary rank is
at least $(D-1)m^2/4$ whereas $h+F_c(A)\le(D+\kappa_D+1)m$.
The all-field matrix upper bound restricts to $\mathcal U$, giving the
matching worst-case product order up to dimensional constants.
The universal rank theorem allows $D=2$; this coefficient application
retains $D\ge3$ throughout.

\section{Component-adapted rectangles and a structured escape}\label{sec:adapted}
We prove Theorem~\ref{thm:adapted} and compare maps on the same coefficient
family and actual coexact domain, without changing either energy.

\subsection{Block spectra and a common finite penalty}
Under the component-block-diagonal hypothesis, a fixed permutation of
gradient coordinates gives actual $D\times D$ blocks $A^k(x)$ and
\[
 Q_A(u)=\sum_{k=1}^h Q_{A^k}(u_k),\qquad
 F_c(A)=\sum_{k=1}^hF_c(A^k).
\]
The spectral identity holds because the full multiset is the disjoint
union of block spectra, including all nonpositive multiplicities and cap
threshold equalities. No spatially varying rotation is performed.
For each actual block apply Theorem~\ref{thm:matrix} with $h=1$.
Here $h=1$ means a single real scalar field with a $D$-coordinate gradient,
not a single spatial direction. It supplies a partition $\mathcal R_k$ with
\[
 \#\mathcal R_k\le a_D+b_DF_c(A^k),\qquad
 V^{-1}Q_{A^k}(u_k)+K_k\sum_{B\in\mathcal R_k}|q_{k,B}(u)|^2
                      \ge cV^{-1}\sum_x|u_k(x)|^2,
\]
where $K_k\ge c$ is finite. All these partitions are fixed from their
blocks before $u$ is chosen. Taking $K=\max_kK_k$ and summing proves
\eqref{eq:adapted-bound}. There is no rescaling of $c$ or $V$ by $h$.
The constants $a_D,b_D$ depend only on $D$.
On the full field space each component's disjoint nonempty block rows are
independent, and different components have independent coordinates.
Thus actual full rank equals $N_{\rm out}$; on a constrained domain it
may decrease. The coefficient-block hypothesis is essential to this proof.

\subsection{Root observations for the strong components}
For $n\ge2$ the unweighted forward torus gradient has first nonzero
Laplacian eigenvalue
\begin{equation}\label{eq:torus-gap}
 \lambda_1=4\sin^2(\pi/n)\ge16/n^2.
\end{equation}
Indeed Fourier symbols give eigenvalues
$\sum_i4\sin^2(\pi k_i/n)$; the least nonzero one has exactly one
nonzero coordinate of smallest distance from zero. The chord bound
$\sin t\ge2t/\pi$ on $[0,\pi/2]$ gives the inequality.
Conjugate pairing gives the same real Poincare bound. The constant
component is handled separately, so no statement about a nonzero gap at
$n=1$ is needed.
For the family \eqref{eq:coexact-family},
\[
 \sigma/c=\alpha_D^{2/D}
 =(3/4)^2[(D/2)(D/(D-2))^{D-1}]^{2/D}\ge9/16.
\]
Consequently $H\lambda_1\ge9crn^{D-2}\ge c$.
Writing $f=\mean_Xf+f_0$ gives
\begin{equation}\label{eq:root-strong}
 V^{-1}H\sum_{x,i}|D_i^+f(x)|^2+c|\mean_Xf|^2
                                      \ge cV^{-1}\sum_x|f(x)|^2.
\end{equation}
The observation here is exactly $q_X(f)=\mean_Xf$.

Only the two real components $(j,\alpha)=(1,1),(2,1)$ are weak.
Apply the $h=1$ matrix theorem to each, obtaining at most
$2a_D+b_D(F_1+F_2)\le2a_D+b_DF_c(A)$ rows in total.
Choose only the root mean for each of the remaining $h-2$ globally strong
components. One maximum of $c$ and the two finite weak-component
penalties yields an all-field inequality with at most
$h-2+2a_D+b_DF_c(A)$ outputs.
Every real component of a coexact field has zero global sum, since it is
a sum of periodic backward differences. Hence all these $h-2$ strong
root rows vanish identically on $\mathcal U$. Restricting the already
proved inequality to that domain gives
\[
 \rank(C_{\rm ad}|_{\mathcal U})\le2a_D+b_DF_c(A).
\]
This does not split a coexact spatial vector into independently variable
spatial components, nor assert independence of the remaining weak rows.
It proves \eqref{eq:adapted-escape} and Theorem~\ref{thm:adapted}.

For this same family $hF_c(A)\ge D\kappa_Dm^2$ and the adapted rank is
at most $2a_D+b_D(\kappa_Dm+1)$. Their ratio tends to zero. In contrast,
Theorem~\ref{thm:sharp-coexact} forces common full-vector maps to have
restricted rank at least $(D-1)m^2/4$. Thus no positive uniform
$\gamma_DhF_c(A)$ lower bound extends to all component-adapted maps.
Refining their different partitions to a common partition and adding
every vector component changes the map and can add the very rank
charged by the common-class lower bound. No optimal adapted rank,
efficient algorithm or arbitrary coupled-matrix additive theorem is inferred.
